The MCU has several \peripheral{GPIOx} peripherals, each of which control up to 8 general purpose input/output (GPIO) pins. Each GPIO peripheral (typically called a port) constitutes a parallel port, allowing all 8 pins in the port to be updated simultaneously. Each GPIO pin may also be configured and updated individually, allowing for design flexibility. Each GPIO pin is referred to as Px.y, where x refers to the port number and y refers to the corresponding bit in the GPIO registers that controls the pin. For example, P2.5 refers to port 2, bit 5.

Every GPIO pin on the MCU can be in one of two modes: primary/GPIO mode, or secondary/peripheral mode. Section \ref{s:pinsConfig} details the secondary function of each GPIO pin.



\subsection{Primary/GPIO Mode Functionality}

In primary/GPIO mode, the GPIO pin state is controlled by the GPIO registers \register{PxDIR}, \register{PxOUT}, and \register{PxREN}. The \register{PxDIR} register determines if the pin is in input or output mode, the \register{PxOUT} register determines if the pin outputs a logic high or logic low level when the pin is in output mode, and the \register{PxREN} register enables or disables the internal pullup resistor when in input mode. Table \ref{t:gpio-config} shows the various states available to every GPIO pin.

\begin{table}[htb]
	\centering
	\begin{tabular}{ c c c c }
	\caption{GPIO Configurations} \label{t:gpio-config}
	\textbf{\register{PxDIR[y]}} & \textbf{\register{PxOUT[y]}} & \textbf{\register{PxREN[y]}} & \textbf{Px.y State} \\ \hline
	0 & - & 0 & High impedance input, pullup resistor disabled \\
	\rowcolor{tablehighlightcolor}
	0 & - & 1 & Pullup resistor enabled, weakly pulled high \\
	1 & 0 & - & Output, strongly driven low \\
	\rowcolor{tablehighlightcolor}
	1 & 1 & - & Output, strongly driven high \\
	\hline
	\end{tabular}
\end{table}

The \register{PxIN} register always reads the digital state of the pin, regardless of if it is in input or output mode, and regardless of if it is in primary/GPIO mode or secondary/peripheral mode. The pin Px.y is a logic low level if \register{PxIN[y]} is a 0, and is a logic high level if it is a 1.

Figure \ref{fig:gpio-pin-diagram} depicts the interaction between the registers and the GPIO pin.

\begin{figure}[htpb]
	\includesvg{figures/gpio-pin-diagram}
	\caption{GPIO Pin Diagram}
	\label{fig:gpio-pin-diagram}
\end{figure}


\subsection{Secondary/Peripheral Mode Functionality}

Most GPIO pins have a secondary/peripheral function which is controlled by one of the peripherals in the MCU. The mode of the GPIO is determined by the \register{PxSEL} register. If \register{PxSEL[y]} is 0, Px.y is in primary/GPIO mode, and if it is 1, Px.y is in secondary/peripheral mode.

In most cases, when in peripheral mode, the peripheral seizes control over the pin's \register{DIR}, \register{OUT}, and \register{REN} controls, overriding \register{PxDIR[y]}, \register{PxOUT[y]}, and \register{PxREN[y]}. For example, when the \pin{SCK0} pin is in peripheral mode, the \peripheral{SPI0} peripheral controls the state of the pin, and the corresponding \register{PxDIR[y]}, \register{PxOUT[y]}, and \register{PxREN[y]} registers have no effect (until \register{PxSEL[y]} becomes 0, of course). However, some pins with secondary/peripheral functions are never seized by the corresponding peripheral, even when in secondary/peripheral mode. In yet others, several but not all of the controls are seized by the peripheral. Considering the \pin{T0CAP0} pin, which is only used by the \peripheral{TIMER0} peripheral as an input, the peripheral does not control the state of the pin at all, and thus it must be manually set as a high impedance input before using it for timing purposes. However, for the \pin{SDA0} pin, the \peripheral{I2C0} peripheral only seizes control over the \register{DIR} and \register{OUT} controls, and allows the \register{REN} control to be manually configurable by \register{PxREN[y]} at all times.



\subsection{Pin Interrupts}

Each GPIO port has a hardware interrupt that can be triggered by any combination of the pins in the port. To enable the interrupt, the \peripheral{GPIOx} interrupt must be unmasked and the desired pin(s) that are allowed to trigger the interrupt must be set as 1s in the \register{PxIE} register.

The interrupt edge select register \register{PxIES} determines which edge will trigger the interrupt on the corresponding pin. The interrupt triggers on a rising edge if \register{PxIES[y]} is 0, and triggers on a falling edge if 1. Note that modifying the value of \register{PxIES[y]} while \register{PxIE[y]} is 1 could immediately generate an interrupt, so the interrupt should be disabled before changing \register{PxIES[y]}.

When the pin Px.y has its interrupt enabled and experiences the edge selected by \register{PxIES[y]}, then the interrupt flag \register{PxIFG[y]} is immediately set to 1 and the \peripheral{GPIOx} ISR is executed. There is one ISR for all of the pins contained in \peripheral{GPIOx}, but the \register{PxIFG[y]} can be used to determine which pin(s) generated the interrupt. Note that \register{PxIFG[y]} must be cleared before the ISR returns, else the ISR will mistakenly re-execute immediately upon the prior return.
